
\FigureLayoutDeclare{img/2020/shanghai/q10_fig1.png}{6.0cm}{8bp 6bp 6bp 8bp}
\FigureLayoutDeclare{img/2020/shanghai/q12_fig1.png}{8.0cm}{9bp 6bp 6bp 7bp}

\chapter{2020年上海卷(秋)}

\section{填空题}

\begin{problem}
已知集合 \(A=\{1,2,4\}\),集合 \(B=\{2,4,5\}\),则 \(A\cap B=\)\fillinblank{}.
\end{problem}

\begin{problem}
计算:\(\lim\limits_{n\to\infty}\frac{n+1}{3n-1}=\)\fillinblank{}.
\end{problem}

\begin{problem}
已知复数 \(z=1-2\i\)(\(\i\) 为虚数单位),则 \(|z|=\)\fillinblank{}.
\end{problem}

\begin{problem}
已知函数 \(f(x)=x^3\),\(f^{-1}(x)\) 是 \(f(x)\) 的反函数,则 \(f^{-1}(x)=\)\fillinblank{}.
\end{problem}

\begin{problem}
已知 \(x\),\(y\) 满足 \(x+y-2\ge0\),\(x+2y-3\le0\),\(y\ge0\),
则 \(z=y-2x\) 的最大值为\fillinblank{}.

\bitmapfigure[max width=6.0cm]{img/2020/shanghai/q5_fig1.png}
\end{problem}

\begin{problem}
已知行列式 \(\left|\begin{matrix}
1 & a & b\\
2 & c & d\\
3 & 0 & 0
\end{matrix}\right|=6\),则 \(\left|\begin{matrix}
a & b\\
c & d
\end{matrix}\right|\)=\fillinblank{}.
\end{problem}

\begin{problem}
已知有四个数 \(1\),\(2\),\(a\),\(b\),这四个数的中位数是 \(3\),平均数是 \(4\),则 \(ab=\)\fillinblank{}.
\end{problem}

\begin{problem}
已知数列 \(\{a_n\}\) 是公差不为零的等差数列,且 \(a_1+a_{10}=a_9\),则
\(\frac{a_1+a_2+\cdots+a_9}{a_{10}}\)=\fillinblank{}.
\end{problem}

\begin{problem}
从 \(6\) 个人挑选 \(4\) 个人去值班,每人值班一天,第一天安排 \(1\) 个人,第二天安排 \(1\) 个人,第三天安排 \(2\) 个人,则共有\fillinblank{} 种安排情况.
\end{problem}

\begin{problem}
已知椭圆 \(C:\frac{x^2}{4}+\frac{y^2}{3}=1\) 的右焦点为 \(F\),直线 \(l\) 经过椭圆右焦点 \(F\),交椭圆 \(C\) 于 \(P\),\(Q\) 两点(点 \(P\) 在第二象限),若点 \(Q\) 关于 \(x\) 轴对称点为 \(Q'\),且满足 \(PQ\perp FQ'\),求直线 \(l\) 的方程是\fillinblank{}.

\bitmapfigure[max width=6.0cm]{img/2020/shanghai/q10_fig1.png}
\end{problem}

\begin{problem}
设 \(a\in\R\),若存在定义域为 \(\R\) 的函数 \(f(x)\) 同时满足下列两个条件:

\begin{enumerate}
\item 对任意的 \(x_0\in\R\),\(f(x_0)\) 的值为 \(x_0\) 或 \(x_0^2\);
\item 关于 \(x\) 的方程 \(f(x)=a\) 无实数解,
\end{enumerate}

则 \(a\) 的取值范围是\fillinblank{}.
\end{problem}

\begin{problem}
已知 \(\bs{a}_1\),\(\bs{a}_2\),\(\bs{b}_1\),\(\bs{b}_2\),\(\ldots\),\(\bs{b}_k\)(\(k\in\N^*\))是平面内两两互不相等的向量,满足 \(|\bs{a}_1-\bs{a}_2|=1\),且 \(|\bs{a}_i-\bs{b}_j|\in\{1,2\}\)(其中 \(i=1\),\(2\),\(j=1\),\(2\),\(\ldots\),\(k\)),则 \(k\) 的最大值是\fillinblank{}.

\bitmapfigure[max width=6.0cm]{img/2020/shanghai/q12_fig1.png}
\end{problem}

\section{单选题}

\begin{problem}
下列不等式恒成立的是(\quad)

\choices
    {\(a^2+b^2\le2ab\)}
    {\(a^2+b^2\ge-2ab\)}
    {\(a+b\ge2\sqrt{|ab|}\)}
    {\(a^2+b^2\le-2ab\)}
\end{problem}

\begin{problem}
已知直线方程 \(3x+4y+1=0\) 的一个参数方程可以是(\quad)

\choices
    {\(x=1+3t\),\(y=-1-4t\)(\(t\) 为参数)}
    {\(x=1-4t\),\(y=-1+3t\)(\(t\) 为参数)}
    {\(x=1-3t\),\(y=-1+4t\)(\(t\) 为参数)}
    {\(x=1+4t\),\(y=1-3t\)(\(t\) 为参数)}
\end{problem}

\begin{problem}
在棱长为 \(10\) 的正方体 \(ABCD-A_1B_1C_1D_1\) 中,\(P\) 为左侧面 \(ADD_1A_1\) 上一点,已知点 \(P\) 到 \(A_1D_1\) 的距离为 \(3\),\(P\) 到 \(AA_1\) 的距离为 \(2\),则过点 \(P\) 且与 \(A_1C\) 平行的直线交正方体于 \(P\),\(Q\) 两点,则 \(Q\) 点所在的平面是(\quad)

\bitmapfigure[max width=6.0cm]{img/2020/shanghai/q15_fig1.png}

\choices
    {\(AA_1B_1B\)}
    {\(BB_1C_1C\)}
    {\(CC_1D_1D\)}
    {\(ABCD\)}
\end{problem}

\begin{problem}
命题 \(p\):存在 \(a\in\R\) 且 \(a\neq0\),对于任意的 \(x\in\R\),使得 \(f(x+a)<f(x)+f(a)\);

命题 \(q_1\):\(f(x)\) 单调递减且 \(f(x)>0\) 恒成立;

命题 \(q_2\):\(f(x)\) 单调递增,存在 \(x_0<0\) 使得 \(f(x_0)=0\),

则下列说法正确的是(\quad)

\choices
    {只有 \(q_1\) 是 \(p\) 的充分条件}
    {只有 \(q_2\) 是 \(p\) 的充分条件}
    {\(q_1\),\(q_2\) 都是 \(p\) 的充分条件}
    {\(q_1\),\(q_2\) 都不是 \(p\) 的充分条件}
\end{problem}

\section{解答题}

\begin{problem}
已知 \(ABCD\) 是边长为 \(1\) 的正方形,正方形 \(ABCD\) 绕 \(AB\) 旋转形成一个圆柱.

\begin{enumerate}
\item 求该圆柱的表面积;
\item 正方形 \(ABCD\) 绕 \(AB\) 逆时针旋转 \(\frac{\pi}{2}\) 至 \(ABC_1D_1\),求线段 \(CD_1\) 与平面 \(ABCD\) 所成的角.
\end{enumerate}

\bitmapfigure[max width=6.0cm]{img/2020/shanghai/q17_fig1.png}
\end{problem}

\begin{problem}
已知函数 \(f(x)=\sin\omega x\),\(\omega>0\).

\begin{enumerate}
\item \(f(x)\) 的周期是 \(4\pi\),求 \(\omega\),并求 \(f(x)=\frac12\) 的解集;
\item 已知 \(\omega=1\),\(g(x)=f^2(x)+\sqrt3\,f(-x)f\!\left(\frac\pi2-x\right)\),\(x\in\left[0,\frac\pi4\right]\),求 \(g(x)\) 的值域.
\end{enumerate}
\end{problem}

\begin{problem}
在研究某市交通情况时,道路密度是指该路段上一定时间内通过的车辆数除以时间,车辆密度是该路段一定时间内通过的车辆数除以该路段的长度,现定义交通流量 \(v=f(x)=\frac{q}{x}\),其中 \(x\) 为道路密度,\(q\) 为车辆密度,且
\(f(x)=100-135\left(\frac13\right)^{\frac{80}{x}}\)(\(0<x<40\)),
\(f(x)=-k(x-40)+85\)(\(40\le x\le80\)).

\begin{enumerate}
\item 若交通流量 \(v>95\),求道路密度 \(x\) 的取值范围;
\item 已知道路密度 \(x=80\) 时,测得交通流量 \(v=50\),求车辆密度 \(q\) 的最大值.
\end{enumerate}
\end{problem}

\begin{problem}
已知双曲线 \(\Gamma_1:\frac{x^2}{4}-\frac{y^2}{b^2}=1\) 与圆 \(\Gamma_2:x^2+y^2=4+b^2\)(\(b>0\))交于点 \(A(x_A,y_A)\)(第一象限),曲线 \(\Gamma\) 为 \(\Gamma_1\),\(\Gamma_2\) 上取满足 \(x>|x_A|\) 的部分.

\begin{enumerate}
\item 若 \(x_A=\sqrt6\),求 \(b\) 的值;
\item 当 \(b=\sqrt5\),\(\Gamma_2\) 与 \(x\) 轴交点记作点 \(F_1\),\(F_2\),\(P\) 是曲线 \(\Gamma\) 上一点,且在第一象限,且 \(|PF_1|=8\),求 \(\angle F_1PF_2\);
\item 过点 \(D\left(0,\frac{b^2}{2}+2\right)\) 斜率为 \(-\frac{b}{2}\) 的直线 \(l\) 与曲线 \(\Gamma\) 只有两个交点,记为 \(M\),\(N\),用 \(b\) 表示 \(\overrightarrow{OM}\cdot\overrightarrow{ON}\),并求 \(\overrightarrow{OM}\cdot\overrightarrow{ON}\) 的取值范围.
\end{enumerate}
\end{problem}

\begin{problem}
已知数列 \(\{a_n\}\) 为有限数列,满足 \(|a_1-a_2|\le|a_1-a_3|\le\cdots\le|a_1-a_m|\),则称 \(\{a_n\}\) 满足性质 \(P\).

\begin{enumerate}
\item 判断数列 \(3\),\(2\),\(5\),\(1\) 和 \(4\),\(3\),\(2\),\(5\),\(1\) 是否具有性质 \(P\),请说明理由;
\item 若 \(a_1=1\),公比为 \(q\) 的等比数列,项数为 \(10\),具有性质 \(P\),求 \(q\) 的取值范围;
\item 若 \(\{a_n\}\) 是 \(1\),\(2\),\(3\),\(\ldots\),\(m\) 的一个排列(\(m\ge4\)),\(\{b_n\}\) 符合 \(b_k=a_{k+1}\)(\(k=1\),\(2\),\(\ldots\),\(m-1\)),\(\{a_n\}\),\(\{b_n\}\) 都具有性质 \(P\),求所有满足条件的数列 \(\{a_n\}\).
\end{enumerate}
\end{problem}

